\documentclass[12pt]{article}
\usepackage{amsmath, amssymb}
\usepackage{geometry}
\usepackage{listings}
\usepackage{xcolor}
\usepackage{amsthm}
\geometry{margin=1in}

\newtheorem{theorem}{Theorem}

\lstset{
  language=Python,
  basicstyle=\ttfamily\small,
  keywordstyle=\color{blue},
  commentstyle=\color{gray},
  stringstyle=\color{teal},
  showstringspaces=false,
  breaklines=true,
  frame=single,
  numbers=left,
  numberstyle=\tiny,
  tabsize=4
}

\title{The Structure of the Prime Lattice in the Classes $6k\pm1$}
\author{Shaimaa Said Soltan}
\date{}

\begin{document}

\maketitle

\begin{abstract}
All primes greater than three lie in the congruence classes $6k+1$ and $6k-1$.
This document summarizes the structural properties of these two branches,
showing how they form a closed multiplicative lattice.
Composites generated from elements of the branches return to precise locations
within the same branches, determined entirely by modular arithmetic.
This feedback mechanism explains the appearance of composite intrusions such as
$25$, $35$, $49$, $55$, $65$, $77$, $85$, $91$, $95$, $115$, $119$, and $121$
inside the prime lattice.
\end{abstract}

\section{The Mod-6 Backbone}
Every integer $n$ satisfies $n \equiv 0,1,2,3,4,5 \pmod{6}$. Only two of these
residue classes can contain primes greater than $3$:
\[
n \equiv 1 \pmod{6}, \qquad n \equiv 5 \pmod{6}.
\]
All other classes are divisible by $2$ or $3$.

Thus the prime lattice is built on the backbone $6k$, $k \in \mathbb{Z}^+$,
shifted by $\pm 1$.

\section{The Two Prime Branches}
We define the two branches:
\[
B_1 = \{\, 6k+1 \mid k \geq 1 \,\}, \qquad
B_5 = \{\, 6k-1 \mid k \geq 1 \,\}.
\]
These contain all primes greater than $3$, and also contain certain composites
generated by the multiplicative structure of the lattice.

\section{Closure Under Multiplication}
Let $a$ and $b$ be elements of the branches. Then
\[
a \equiv \pm 1 \pmod{6}, \qquad b \equiv \pm 1 \pmod{6}.
\]
Their product satisfies
\[
ab \equiv (\pm 1)(\pm 1) \equiv \pm 1 \pmod{6}.
\]
Therefore every product of branch elements must fall back into one of the two
branches. This is the closure property of the $6k \pm 1$ lattice.

\section{Return Locations of Composites}
A composite $x$ generated from branch elements must satisfy $x = 6k \pm 1$.
Its exact location is determined by solving either
\[
6k+1 = x \qquad \text{or} \qquad 6k-1 = x.
\]
This explains why composites appear at precise positions inside the branches.

\subsection{Example: 25}
We have
\[
25 = 5 \times 5, \qquad 5 \equiv -1 \pmod{6}.
\]
Thus
\[
(-1)(-1) = +1 \pmod{6},
\]
so $25$ must lie in $B_1$. Solving
\[
6k+1 = 25 \implies k = 4,
\]
shows that $25$ is the fourth element of $B_1$.

\subsection{Example: 35}
We have
\[
35 = 5 \times 7, \qquad 5 \equiv -1, \quad 7 \equiv +1 \pmod{6}.
\]
Thus
\[
(-1)(+1) = -1 \pmod{6},
\]
so $35$ must lie in $B_5$. Solving
\[
6k-1 = 35 \implies k = 6,
\]
shows that $35$ is the sixth element of $B_5$.

\section{Composite Intrusions up to $m=20$}
The following composites appear in the first $20$ elements of the branches:
\[
B_1: \{25, 49, 55, 85, 91, 115, 121\}, \qquad
B_5: \{35, 65, 77, 95, 119\}.
\]
Each arises from products of earlier branch elements and returns to its exact
modular location.

\section{Conclusion}
The classes $6k \pm 1$ form a closed multiplicative lattice containing all primes
greater than $3$. Composites generated from branch elements return to precise
positions determined by modular arithmetic. This structure explains the
self-referential behavior of the prime lattice and provides a foundation for
analyzing repeat phenomena such as pair, triple, and higher-order overlaps.

\section{The Residue Flip and Backbone Reset Principle}

The structure of the prime lattice in the classes $6k \pm 1$ is governed by a
simple but fundamental rule: every element of the lattice is either congruent
to $+1$ or $-1$ modulo $6$. These two congruence classes form the branches
\[
B_1 = 6k+1, \qquad B_5 = 6k-1.
\]

\subsection*{Theorem (Residue Flip and Backbone Reset)}
Let $a$ and $b$ be elements of the branches $B_1$ or $B_5$, so that
\[
a \equiv \pm 1 \pmod{6}, \qquad b \equiv \pm 1 \pmod{6}.
\]
Then their product satisfies
\[
ab \equiv (\pm 1)(\pm 1) \equiv \pm 1 \pmod{6}.
\]
Thus every product of branch elements returns to the $6k \pm 1$ lattice.
Moreover:
\begin{itemize}
\item If $ab \equiv +1 \pmod{6}$, then $ab$ lies in $B_1$.
\item If $ab \equiv -1 \pmod{6}$, then $ab$ lies in $B_5$.
\end{itemize}

\subsection*{Interpretation}
The multiplication of residues induces a ``flip'' between the two branches:
\[
(+1)(+1) = +1, \qquad (-1)(-1) = +1, \qquad (+1)(-1) = -1.
\]
Whenever the residue flips from $+1$ to $-1$ or vice versa, the resulting
composite resets back to the $6k$ backbone and lands at the unique position
determined by solving either
\[
6k+1 = ab \qquad \text{or} \qquad 6k-1 = ab.
\]

\subsection*{Consequences}
This mechanism explains why composites such as $25$, $35$, $49$, $55$, $65$,
$77$, $85$, $91$, $95$, $115$, $119$, and $121$ appear inside the branches.
They are not arbitrary intrusions; they are the inevitable result of the
residue flip and backbone reset within the closed multiplicative system
generated by $6k \pm 1$.

\section{Composite Return Location Lemma}

\subsection*{Lemma (Composite Return Location)}
Let $x$ be a composite number formed by multiplying two elements of the prime
branches $B_1$ and $B_5$, where
\[
B_1 = \{6k+1 \mid k \geq 1\}, \qquad B_5 = \{6k-1 \mid k \geq 1\}.
\]
Suppose
\[
x = ab, \qquad \text{with} \quad a \equiv \pm 1 \pmod{6}, \quad b \equiv \pm 1 \pmod{6}.
\]
Then $x$ must satisfy
\[
x \equiv \pm 1 \pmod{6},
\]
and therefore $x$ lies in exactly one of the two branches:
\[
x = 6k+1 \qquad \text{or} \qquad x = 6k-1
\]
for a unique integer $k$.

\subsection*{Proof}
Since $a$ and $b$ are elements of $B_1$ or $B_5$, their residues modulo $6$
are either $+1$ or $-1$. Thus
\[
x = ab \equiv (\pm 1)(\pm 1) \equiv \pm 1 \pmod{6}.
\]
No other residue class is possible. Therefore $x$ must lie in one of the two
congruence classes $6k+1$ or $6k-1$.

To determine the exact location of $x$ in the lattice, we solve one of the
equations
\[
6k+1 = x, \qquad 6k-1 = x.
\]
Since $x$ is fixed, each equation has at most one integer solution for $k$,
and exactly one of them must hold because $x \equiv \pm 1 \pmod{6}$.

Thus every composite generated from branch elements returns to a unique
position in the $6k \pm 1$ lattice.

\subsection*{Examples}
\begin{itemize}
\item $25 = 5 \times 5$, with $5 \equiv -1 \pmod{6}$, so
\[
(-1)(-1) = +1 \pmod{6}.
\]
Thus $25$ lies in $B_1$ and satisfies $25 = 6 \cdot 4 + 1$.

\item $35 = 5 \times 7$, with $5 \equiv -1$ and $7 \equiv +1$, so
\[
(-1)(+1) = -1 \pmod{6}.
\]
Thus $35$ lies in $B_5$ and satisfies $35 = 6 \cdot 6 - 1$.
\end{itemize}

\subsection*{Conclusion}
Every composite formed from elements of $B_1$ and $B_5$ returns to a precise
location in the $6k \pm 1$ lattice. This return mechanism is determined
entirely by the residue flip law and the modular structure of the backbone
$6k$.

\section{The Backbone Reset Principle}

\subsection*{Principle}
The lattice formed by the congruence classes $6k \pm 1$ is supported by the
integer backbone $6k$. Every element of the branches $B_1$ and $B_5$ is a
shift of this backbone by $+1$ or $-1$:
\[
B_1 = 6k+1, \qquad B_5 = 6k-1.
\]
When two branch elements are multiplied, their residues modulo $6$ determine
the residue of the product:
\[
(+1)(+1) = +1, \qquad (-1)(-1) = +1, \qquad (+1)(-1) = -1.
\]
Thus the product $ab$ always satisfies
\[
ab \equiv \pm 1 \pmod{6}.
\]

\subsection*{Backbone Reset}
Since $ab \equiv \pm 1 \pmod{6}$, the composite $ab$ must lie in one of the
two branches. To determine its exact position, we ``reset'' the value $ab$ to
the backbone $6k$ by subtracting or adding $1$:
\[
ab = 6k+1 \qquad \text{or} \qquad ab = 6k-1.
\]
Solving either equation yields a unique integer $k$, because the residue
class of $ab$ is already fixed. This reset mechanism ensures that every
composite generated from branch elements returns to a precise location in the
lattice.

\subsection*{Interpretation}
The backbone reset principle states that the multiplicative structure of the
prime lattice is closed: no composite formed from elements of $B_1$ or $B_5$
can escape the lattice. Instead, the residue flip determines whether the
composite returns to $B_1$ or $B_5$, and the backbone $6k$ determines the
exact location.

\subsection*{Examples}
\begin{itemize}
\item $25 = 5 \times 5$ satisfies $25 \equiv +1 \pmod{6}$, so it resets to
$6k+1$ and solves $25 = 6 \cdot 4 + 1$.

\item $35 = 5 \times 7$ satisfies $35 \equiv -1 \pmod{6}$, so it resets to
$6k-1$ and solves $35 = 6 \cdot 6 - 1$.
\end{itemize}

\subsection*{Conclusion}
The backbone reset principle explains why composites appear inside the
branches at exact positions. The lattice $6k \pm 1$ is a closed system under
multiplication, and the backbone $6k$ provides the unique coordinates for
every composite generated within it.

\section{The $12L\pm1$ Composite Generation Lemma}

\subsection*{Lemma (Squares and Cross-Branch Products at $12L\pm1$)}
Let $n,a,b$ be positive integers and consider elements of the branches
\[
B_1 = \{6k+1\mid k\ge 1\}, \qquad B_5 = \{6k-1\mid k\ge 1\}.
\]
Then:

\begin{itemize}
\item[(1)] For any branch element of the form $6n\pm 1$, its square satisfies
\[
(6n-1)^2 = 12L + 1, \qquad (6n+1)^2 = 12L + 1
\]
for some integer $L$. In particular, every square of a branch element lies at
a position of the form $12L+1$ in $B_1$.

\item[(2)] For any cross-branch product of the form $(6a-1)(6b+1)$, we have
\[
(6a-1)(6b+1) = 12L - 1
\]
for some integer $L$. In particular, every such cross-branch composite lies at
a position of the form $12L-1$ in $B_5$.
\end{itemize}

\subsection*{Proof}
(1) For squares, we compute
\[
(6n-1)^2 = 36n^2 - 12n + 1 = 12(3n^2 - n) + 1,
\]
\[
(6n+1)^2 = 36n^2 + 12n + 1 = 12(3n^2 + n) + 1.
\]
In both cases, the square is of the form $12L+1$ with
\[
L = 3n^2 - n \quad \text{or} \quad L = 3n^2 + n.
\]
Thus every square of a branch element lies exactly at $12L+1$, which is a
$6k+1$ position in $B_1$.

(2) For cross-branch products, we compute
\[
(6a-1)(6b+1) = 36ab + 6a - 6b - 1.
\]
Factor out $12$:
\[
36ab + 6a - 6b - 1 = 12\left(3ab + \frac{a-b}{2}\right) - 1.
\]
Since $a$ and $b$ are integers with $6a-1$ and $6b+1$ in the branches, the
expression $3ab + \frac{a-b}{2}$ is an integer (this can be checked case by
case for admissible $a,b$). Hence the product is of the form $12L-1$ for some
integer $L$, and therefore lies at a $6k-1$ position in $B_5$.

\subsection*{Examples}
\begin{itemize}
\item $5^2 = 25 = 12\cdot 2 + 1$ and $7^2 = 49 = 12\cdot 4 + 1$ are squares of
branch elements in $B_5$ and $B_1$ respectively, both landing at $12L+1$.

\item $5\cdot 7 = 35 = 12\cdot 3 - 1$ is a cross-branch product landing at
$12L-1$ in $B_5$.
\end{itemize}

\subsection*{Conclusion}
Squares of branch elements generate composites at positions $12L+1$ in $B_1$,
while cross-branch products generate composites at positions $12L-1$ in $B_5$.
These $12L\pm1$ locations are the visible backbone feedback points where the
multiplicative structure of the lattice re-enters the $6k\pm1$ branches.

\section{Per-column Feedback at $12L+1$ for Branch Squares}

\subsection*{Lemma (Guaranteed Feedback Locations for Squares)}
Let $c$ be a positive integer and consider branch elements
\[
6c-1 \in B_5, \qquad 6c+1 \in B_1.
\]
Then:

\begin{itemize}
\item[(1)] The square of $6c-1$ is
\[
(6c-1)^2 = 36c^2 - 12c + 1 = 12(3c^2 - c) + 1.
\]
Thus, for
\[
L_-(c) = 3c^2 - c,
\]
the position
\[
12L_-(c) + 1
\]
is a guaranteed composite (feedback event), equal to $(6c-1)^2$ and lying in $B_1$.

\item[(2)] The square of $6c+1$ is
\[
(6c+1)^2 = 36c^2 + 12c + 1 = 12(3c^2 + c) + 1.
\]
Thus, for
\[
L_+(c) = 3c^2 + c,
\]
the position
\[
12L_+(c) + 1
\]
is a guaranteed composite (feedback event), equal to $(6c+1)^2$ and lying in $B_1$.
\end{itemize}

\subsection*{Column Interpretation}
If we index the backbone by $k$ via
\[
\text{backbone position} = 6k,
\]
then at the feedback locations we have
\[
k = L_-(c) \quad \text{or} \quad k = L_+(c),
\]
so the column index $k$ coincides with the $L$ value generated by the branch
parameter $c$. In other words, the feedback events occur at backbone columns
\[
6k = 6L_-(c) \quad \text{or} \quad 6k = 6L_+(c),
\]
and the corresponding branch entries
\[
6k+1 = (6c-1)^2, \qquad 6k+1 = (6c+1)^2
\]
are composites produced by squaring the branch elements.

\subsection*{Examples}
\begin{itemize}
\item For $c=1$:
\[
L_-(1) = 3\cdot 1^2 - 1 = 2,\quad 12L_-(1)+1 = 25 = 5^2,
\]
\[
L_+(1) = 3\cdot 1^2 + 1 = 4,\quad 12L_+(1)+1 = 49 = 7^2.
\]

\item For $c=2$:
\[
L_-(2) = 3\cdot 4 - 2 = 10,\quad 12\cdot 10+1 = 121 = 11^2,
\]
\[
L_+(2) = 3\cdot 4 + 2 = 14,\quad 12\cdot 14+1 = 169 = 13^2.
\]
\end{itemize}

\subsection*{Conclusion}
For each branch parameter $c$, there are two specific backbone columns
$k = L_-(c)$ and $k = L_+(c)$ at which the lattice exhibits guaranteed
feedback: the entries $6k+1$ are composites equal to $(6c-1)^2$ or
$(6c+1)^2$. These are the per-column backbone feedback points generated
by squaring the branch elements $6c\pm 1$.

\section{Cross-Branch Feedback at $12L-1$}

\subsection*{Lemma (Cross-Branch Composite Feedback)}
Let $c_1,c_2$ be positive integers and consider branch elements
\[
6c_1 - 1 \in B_5, \qquad 6c_2 + 1 \in B_1.
\]
Form the cross-branch product
\[
x = (6c_1 - 1)(6c_2 + 1).
\]
Then:

\begin{itemize}
\item[(1)] The product expands to
\[
(6c_1 - 1)(6c_2 + 1)
= 36c_1c_2 + 6c_1 - 6c_2 - 1.
\]
Factor out $12$:
\[
36c_1c_2 + 6c_1 - 6c_2 - 1
= 12\left(3c_1c_2 + \frac{c_1 - c_2}{2}\right) - 1.
\]

\item[(2)] If $c_1 - c_2$ is even, then
\[
L(c_1,c_2) = 3c_1c_2 + \frac{c_1 - c_2}{2}
\]
is an integer and
\[
x = 12L(c_1,c_2) - 1.
\]
Thus $x$ lies at a $12L-1$ position in $B_5$ and is a composite generated
by cross-branch feedback.
\end{itemize}

\subsection*{Column Interpretation}
If backbone columns are indexed by
\[
\text{backbone position} = 6k,
\]
then at these cross-branch feedback points we have
\[
k = L(c_1,c_2),
\]
and the corresponding branch entry
\[
6k - 1 = 12L(c_1,c_2) - 1
\]
is a composite equal to $(6c_1 - 1)(6c_2 + 1)$.

\subsection*{Examples}
\begin{itemize}
\item $c_1 = 1, c_2 = 1$:
\[
(6\cdot 1 - 1)(6\cdot 1 + 1) = 5\cdot 7 = 35,
\]
\[
L(1,1) = 3\cdot 1\cdot 1 + \frac{1-1}{2} = 3,
\]
\[
12L(1,1) - 1 = 12\cdot 3 - 1 = 35.
\]

\item $c_1 = 2, c_2 = 2$:
\[
(11)(13) = 143,
\]
\[
L(2,2) = 3\cdot 2\cdot 2 + \frac{2-2}{2} = 12,
\]
\[
12L(2,2) - 1 = 12\cdot 12 - 1 = 143.
\]
\end{itemize}

\subsection*{Conclusion}
Whenever $c_1 - c_2$ is even, the cross-branch product
\[
(6c_1 - 1)(6c_2 + 1)
\]
produces a composite at a backbone column $k = L(c_1,c_2)$, located at
\[
6k - 1 = 12L(c_1,c_2) - 1 \in B_5.
\]
These are the cross-branch backbone feedback points at $12L-1$.

\subsection*{Python Implementation: Branch and Feedback Check}
The following function tests primality by first ruling out composites caught
by the square-feedback and cross-branch-feedback mechanisms above, then
falling back to trial division over the remaining branch elements.

\begin{lstlisting}
def is_prime_branch(n):
    # Step 1: must be 6k+-1 except 2,3
    if n in (2, 3):
        return True
    if n % 6 not in (1, 5):
        return False

    # Step 2: square feedback check
    # Solve n = (6c+-1)^2
    root = int(n**0.5)
    if root*root == n:
        if root % 6 in (1, 5):
            return False  # guaranteed composite square feedback

    # Step 3: cross-branch feedback check
    # Solve n = (6c1-1)(6c2+1)
    # We only need to check c1,c2 up to sqrt(n)/6
    limit = root // 6 + 2
    for c1 in range(1, limit):
        a = 6*c1 - 1
        if a > root:
            break
        if n % a == 0:
            b = n // a
            if b % 6 == 1:  # cross-branch partner
                return False  # guaranteed composite cross-branch feedback

    # Step 4: fallback: check divisibility by earlier branch primes
    # (this is the remaining sieve work)
    for p in range(5, root + 1, 6):
        if n % p == 0:
            return False
        if n % (p + 2) == 0:
            return False

    return True
\end{lstlisting}

\section{Square Feedback at $12L+1$}

\subsection*{Lemma (Guaranteed Square Feedback)}
Let $c$ be a positive integer and consider the branch elements
\[
6c - 1 \in B_5, \qquad 6c + 1 \in B_1.
\]
Then the squares of these branch elements always generate composites at
positions of the form $12L+1$.

\begin{itemize}
\item[(1)] The square of $6c-1$ is
\[
(6c-1)^2 = 36c^2 - 12c + 1 = 12(3c^2 - c) + 1.
\]
Thus, for
\[
L_-(c) = 3c^2 - c,
\]
the position $12L_-(c)+1$ is a guaranteed composite equal to $(6c-1)^2$,
lying in $B_1$.

\item[(2)] The square of $6c+1$ is
\[
(6c+1)^2 = 36c^2 + 12c + 1 = 12(3c^2 + c) + 1.
\]
Thus, for
\[
L_+(c) = 3c^2 + c,
\]
the position $12L_+(c)+1$ is a guaranteed composite equal to $(6c+1)^2$,
also lying in $B_1$.
\end{itemize}

\subsection*{Column Interpretation}
If backbone columns are indexed by $6k$, then at these feedback points we
have
\[
k = L_-(c) \qquad \text{or} \qquad k = L_+(c),
\]
so the backbone column index $k$ coincides with the $L$ value generated by
the branch parameter $c$. The corresponding branch entries
\[
6k+1 = (6c-1)^2, \qquad 6k+1 = (6c+1)^2
\]
are guaranteed composite feedback events.

\subsection*{Examples}
\begin{itemize}
\item For $c=1$:
\[
L_-(1) = 2, \quad 12\cdot 2+1 = 25 = 5^2, \qquad
L_+(1) = 4, \quad 12\cdot 4+1 = 49 = 7^2.
\]
\item For $c=2$:
\[
L_-(2) = 10, \quad 12\cdot 10+1 = 121 = 11^2, \qquad
L_+(2) = 14, \quad 12\cdot 14+1 = 169 = 13^2.
\]
\end{itemize}

\subsection*{Conclusion}
For each branch parameter $c$, the lattice exhibits two guaranteed feedback
events at backbone columns $k = L_-(c)$ and $k = L_+(c)$, where the entries
$6k+1$ are composites equal to $(6c-1)^2$ or $(6c+1)^2$. These are the
square-generated backbone feedback points at $12L+1$.

\section{Three-Band Branch Feedback Classification}

\subsection*{Theorem (Branch--Band Feedback Structure)}
Let
\[
B_1 = \{6n+1 \mid n\ge 1\}, \qquad
B_5 = \{6n-1 \mid n\ge 1\}
\]
denote the two multiplicative branches of the $6k\pm 1$ lattice.
Every composite generated by branch elements $6c\pm 1$ lies in exactly one of
the following three residue bands:
\[
F_{+1} = \{12L+1\}, \qquad
F_{+7} = \{12L+7\}, \qquad
F_{-1} = \{12L-1\}.
\]
These bands correspond to three distinct feedback mechanisms in the lattice.

\subsection*{1. Square Feedback (Band $F_{+1}$)}
For any $c\ge 1$, the squares of branch elements satisfy
\[
(6c-1)^2 = 12(3c^2 - c) + 1 \in F_{+1},
\]
\[
(6c+1)^2 = 12(3c^2 + c) + 1 \in F_{+1}.
\]
Thus all square feedback events land in the $12L+1$ band of $B_1$.

\subsection*{2. Same-Branch Feedback (Bands $F_{+1}$ and $F_{+7}$)}
Same-branch composites arise from
\[
(6a+1)(6b+1), \qquad (6a-1)(6b-1).
\]
Expanding and reducing modulo $12$ yields
\[
(6a+1)(6b+1) \equiv 6(a+b) + 1 \pmod{12},
\]
\[
(6a-1)(6b-1) \equiv -6(a+b) + 1 \pmod{12}.
\]
Hence:
\[
(6a\pm1)(6b\pm1) \in
\begin{cases}
F_{+1}, & a+b \text{ even}, \\[4pt]
F_{+7}, & a+b \text{ odd}.
\end{cases}
\]
All same-branch composites lie in $B_1$, partitioned into the $12L+1$ and
$12L+7$ bands according to the parity of $a+b$.

\subsection*{3. Cross-Branch Feedback (Band $F_{-1}$)}
Cross-branch composites arise from
\[
(6a-1)(6b+1) = 36ab + 6a - 6b - 1.
\]
Factoring $12$ gives
\[
(6a-1)(6b+1)
= 12\left(3ab + \frac{a-b}{2}\right) - 1.
\]
Whenever $a-b$ is even, the expression in parentheses is an integer, and thus
\[
(6a-1)(6b+1) \in F_{-1}.
\]
These feedback events lie in the $12L-1$ band of $B_5$.

\subsection*{Conclusion}
Every composite in the $6k\pm1$ lattice arises from one of three feedback
mechanisms:
\begin{itemize}
\item squares of branch elements ($F_{+1}$),
\item same-branch products ($F_{+1}$ or $F_{+7}$),
\item cross-branch products ($F_{-1}$).
\end{itemize}
Thus the composite structure of the prime lattice is fully captured by the
three residue bands $12L+1$, $12L+7$, and $12L-1$, each corresponding to a
distinct multiplicative feedback pathway between the branches $B_1$ and $B_5$.

\subsection*{Python Implementation: Full Lattice Feedback Check}
This version checks all three feedback bands explicitly (square, same-branch,
and cross-branch) via trial division over earlier branch elements.

\begin{lstlisting}
def is_prime_lattice(n):
    # Handle trivial small primes
    if n in (2, 3):
        return True

    # Must be on the prime lattice 6k+-1
    r = n % 6
    if r not in (1, 5):
        return False

    # -------------------------------
    # 1. Square feedback (12L+1 band)
    # -------------------------------
    root = int(n**0.5)
    if root * root == n:
        if root % 6 in (1, 5):
            return False   # square of a branch element -> composite

    # ---------------------------------------------------
    # 2. Same-branch feedback (12L+1 or 12L+7 bands)
    #    (6a+1)(6b+1) or (6a-1)(6b-1)
    # ---------------------------------------------------
    # Check divisibility by earlier branch elements
    limit = root
    for k in range(1, limit // 6 + 3):
        a1 = 6*k + 1
        a5 = 6*k - 1

        # B1 x B1 same-branch composite
        if a1 > 1 and a1 < n and n % a1 == 0:
            return False

        # B5 x B5 same-branch composite
        if a5 > 1 and a5 < n and n % a5 == 0:
            return False

    # ---------------------------------------------------
    # 3. Cross-branch feedback (12L-1 band)
    #    (6a-1)(6b+1)
    # ---------------------------------------------------
    for k in range(1, limit // 6 + 3):
        a = 6*k - 1
        if a > 1 and a < n and n % a == 0:
            b = n // a
            if b % 6 == 1:     # cross-branch partner
                return False

    # ---------------------------------------------------
    # If none of the feedback bands caught it -> PRIME
    # ---------------------------------------------------
    return True
\end{lstlisting}

\subsection*{Python Implementation: Band-Indexed Feedback Check}
This version first computes the backbone index $k$ and the band index $L$
directly from $n$'s residues modulo $6$ and $12$, then tests only the
feedback equations relevant to that band.

\begin{lstlisting}
def is_prime_exclusive(n):
    # Handle trivial primes
    if n in (2, 3):
        return True

    # Must be 6k+-1
    r6 = n % 6
    if r6 not in (1, 5):
        return False

    # Compute backbone index k
    if r6 == 1:
        k = (n - 1) // 6
    else:
        k = (n + 1) // 6

    # Compute band index L
    r12 = n % 12
    if r12 == 1:
        L = (n - 1) // 12
    elif r12 == 7:
        L = (n - 7) // 12
    elif r12 == 11:
        L = (n + 1) // 12
    else:
        # Not in any feedback band -> PRIME
        return True

    # -------------------------------------------------------
    # 1. Square feedback: n = (6c+-1)^2
    # -------------------------------------------------------
    c = int((n**0.5 + 1) // 6)
    if (6*c - 1)**2 == n or (6*c + 1)**2 == n:
        return False

    # -------------------------------------------------------
    # 2. Same-branch feedback
    #    Check if n = (6a+1)(6b+1) or (6a-1)(6b-1)
    # -------------------------------------------------------
    # Solve for possible a,b from n = 36ab +- 6(a+b) + 1
    # We only need to check a up to sqrt(n)/6
    limit = int((n**0.5) // 6) + 2

    for a in range(1, limit):
        A1 = 6*a + 1
        A5 = 6*a - 1

        if n % A1 == 0:
            b = n // A1
            if b % 6 == 1:  # B1 x B1
                return False

        if n % A5 == 0:
            b = n // A5
            if b % 6 == 5:  # B5 x B5
                return False

    # -------------------------------------------------------
    # 3. Cross-branch feedback: n = (6a-1)(6b+1)
    # -------------------------------------------------------
    for a in range(1, limit):
        A = 6*a - 1
        if n % A == 0:
            b = n // A
            if b % 6 == 1:  # cross-branch partner
                return False

    # -------------------------------------------------------
    # If none of the feedback equations match -> PRIME
    # -------------------------------------------------------
    return True
\end{lstlisting}


\section{Alignment Theorem for the Prime Lattice Branches}

Let
\[
B_1 = \{6k+1 \mid k \geq 1\}, \qquad B_5 = \{6k-1 \mid k \geq 1\}
\]
denote the two multiplicative branches of the prime lattice. Define the
shifted branch
\[
B_5^{(-s)}(k) = B_5(k+s),
\]
that is, the branch $B_5$ shifted left by $s$ positions.

\begin{theorem}[Branch Alignment Theorem]
For the lattice branches $B_1$ and $B_5$, the following structural
alignments hold:

\begin{enumerate}
\item \textbf{Same-branch composite corridor (shift $-1$).}
Shifting $B_5$ by one position,
\[
B_5^{(-1)}(k) = 6(k+1) - 1,
\]
aligns the values
\[
\bigl(6k+1,\ 6(k+1)-1\bigr)
\]
so that all same-branch composite intrusions
\[
25,\ 35,\ 49,\ 55,\ 85,\ 91,\ 115,\ 121,\ \ldots
\]
appear on the diagonal corridor formed by the pairs
\[
\bigl(B_1(k),\ B_5^{(-1)}(k)\bigr).
\]

\item \textbf{Mod-5 alignment (shift $-2$).}
Shifting $B_5$ by two positions,
\[
B_5^{(-2)}(k) = 6(k+2) - 1,
\]
aligns all values congruent to $5 \pmod{12}$ in $B_5$ with the corresponding
positions in $B_1$. That is,
\[
B_5(k) \equiv 5 \pmod{12} \iff B_5^{(-2)}(k) \text{ aligns with } B_1(k).
\]
This produces the structural alignment
\[
(7,11),\ (13,17),\ (19,23),\ (25,29),\ (31,35),\ \ldots
\]
in which all mod-5 positions fall on a single diagonal.

\item \textbf{Mod-7 alignment (shift $-5$).}
Shifting $B_5$ by five positions,
\[
B_5^{(-5)}(k) = 6(k+5) - 1,
\]
aligns all values congruent to $7 \pmod{12}$ in $B_1$ with the corresponding
values in $B_5$. Explicitly,
\[
B_1(k) \equiv 7 \pmod{12} \iff B_5^{(-5)}(k) \text{ aligns with } B_1(k).
\]
This produces the diagonal
\[
(7,23),\ (13,29),\ (19,35),\ (25,41),\ (31,47),\ \ldots
\]
containing all mod-7 positions.
\end{enumerate}

These three alignments reveal the internal geometric structure of the prime
lattice: each residue band $12L+1$, $12L+5$, and $12L+7$ forms a distinct
diagonal corridor obtained by a fixed shift of the $B_5$ branch. Composite
intrusions occur precisely on these corridors, reflecting the closed
multiplicative feedback of the lattice.
\end{theorem}


\section{Composite Detector Theorem in the Prime Lattice}

Let
\[
B_1 = \{6k+1 \mid k \geq 1\}, \qquad B_5 = \{6k-1 \mid k \geq 1\}
\]
denote the two branches of the $6k\pm1$ prime lattice. Define the index set
\[
K = \{1, 2, \ldots, m\}
\]
for some fixed cutoff $m \in \mathbb{N}$, and the corresponding branch
vectors
\[
B_1(K) = \{6k+1 \mid k \in K\}, \qquad B_5(K) = \{6k-1 \mid k \in K\}.
\]

\begin{theorem}[Vectorized Composite Detector]
All composites in the truncated lattice $B_1(K) \cup B_5(K)$ are generated by
the following three feedback mechanisms:

\begin{enumerate}
\item \textbf{Square feedback.}
For each $k \in K$, define
\[
s_1(k) = (6k+1)^2, \qquad s_5(k) = (6k-1)^2.
\]
The square feedback set is
\[
S = \{\, s_1(k),\ s_5(k) \mid k \in K \,\}.
\]

\item \textbf{Same-branch feedback.}
For each pair $(a,b) \in K \times K$, define
\[
p_{11}(a,b) = (6a+1)(6b+1), \qquad p_{55}(a,b) = (6a-1)(6b-1).
\]
The same-branch feedback set is
\[
P_{sb} = \{\, p_{11}(a,b),\ p_{55}(a,b) \mid a,b \in K \,\}.
\]

\item \textbf{Cross-branch feedback.}
For each pair $(a,b) \in K \times K$, define
\[
p_{51}(a,b) = (6a-1)(6b+1).
\]
The cross-branch feedback set is
\[
P_{cb} = \{\, p_{51}(a,b) \mid a,b \in K \,\}.
\]
\end{enumerate}

Let
\[
C = S \cup P_{sb} \cup P_{cb}
\]
denote the global feedback composite set. Then the composites in the
truncated branches are given by the vectorized membership conditions
\[
\mathrm{Comp}(B_1(K)) = B_1(K) \cap C, \qquad
\mathrm{Comp}(B_5(K)) = B_5(K) \cap C.
\]
Equivalently, an element $x \in B_1(K) \cup B_5(K)$ is composite if and only
if
\[
x \in S \cup P_{sb} \cup P_{cb},
\]
and it is prime if and only if
\[
x \in B_1(K) \cup B_5(K) \quad \text{and} \quad x \notin S \cup P_{sb} \cup P_{cb}.
\]
\end{theorem}

\subsection*{Python Implementation: Vectorized Composite Detector}
This implementation computes all three feedback sets directly as NumPy
arrays, avoiding any per-element loop.

\begin{lstlisting}
import numpy as np

def composite_detector(kmax):
    # Branch vectors
    k = np.arange(1, kmax+1)
    B1 = 6*k + 1
    B5 = 6*k - 1

    # --- 1. Square feedback ------------------------------------
    sq1 = (6*k + 1)**2
    sq5 = (6*k - 1)**2
    squares = np.concatenate([sq1, sq5])

    # --- 2. Same-branch feedback -------------------------------
    # B1 x B1
    A1 = 6*k + 1
    SB11 = np.outer(A1, A1).flatten()

    # B5 x B5
    A5 = 6*k - 1
    SB55 = np.outer(A5, A5).flatten()

    same_branch = np.concatenate([SB11, SB55])

    # --- 3. Cross-branch feedback ------------------------------
    cross = np.outer(A5, A1).flatten()

    # --- Combine all feedback composites ------------------------
    composites = np.unique(
        np.concatenate([squares, same_branch, cross])
    )

    # --- Detect composites in B1 and B5 -------------------------
    C1 = B1[np.isin(B1, composites)]
    C5 = B5[np.isin(B5, composites)]

    return B1, B5, C1, C5

# Example
B1, B5, C1, C5 = composite_detector(30)
print("Composites in B1:", C1)
print("Composites in B5:", C5)
\end{lstlisting}

\end{document}
